\section{Methode 3: Risikomonitoring und Risikoreporting} \label{sec:risiko_monitoring}

Nach Identifikation, Bewertung und Maßnahmenplanung müssen Risiken fortlaufend
überwacht werden: Annahmen ändern sich, Abhängigkeiten entstehen, Risiken eskalieren.
Risikomonitoring und -reporting zielen auf frühe Warnsignale und nachvollziehbare
Entscheidungen gegenüber Stakeholdern \cite{Tiemeyer2023}.
Für Data-Science-Projekte umfasst Monitoring neben Zeit/Budget/Scope auch zentrale
Artefakte (Daten, Labels, Feature-Pipelines, Experimente, Modellversionen, Modellleistung im Betrieb),
da dort spezifische, sonst schwer sichtbare Risiken entstehen \cite{Sculley2015,Amershi2019}.

\subsection{Ziel, Taktung und Verantwortlichkeiten}
Ein wirksames Monitoring benötigt eine klare Taktung und eindeutige Zuständigkeiten.
Als Minimalset hat sich ein wiederkehrender Zyklus bewährt. Typische Elemente sind:
\begin{itemize}
	\item \textbf{Regeltermine für Risikoreviews:}
				kurze, feste Checkpoints zur Erfassung neuer Risiken und Statusänderungen.
	\item \textbf{Aktualisierung des Risikoregisters:}
				Pflege von Bewertung, Wirksamkeit bestehender Kontrollen und offenen Aufgaben.
	\item \textbf{Verantwortungszuordnung:}
				Zuordnung von Umsetzung und Freigabe je priorisiertem Risiko.
	\item \textbf{Review von Datenindikatoren und Modellindikatoren:}
				Kurzprüfung von Daten-/Labelqualität, Modellgüte sowie Drift-Anzeichen.
	\item \textbf{Nachvollziehbarkeit der Arbeitsergebnisse:}
				Entscheidungen auf reproduzierbaren Artefakten (Versionen, Experimente, Freigaben).
\end{itemize}

\subsection{Frühwarnmechanismen und Eskalation}
Monitoring ist besonders effektiv,
wenn es neben der reinen Statuspflege auch messbare Warnsignale nutzt.
Solche Frühwarnindikatoren (Key Risk Indicators) operationalisieren Risiken über
beobachtbare Größen \cite{Tiemeyer2023}.

Wesentlich ist die Definition von Schwellenwerten und Reaktionen:
\begin{itemize}
  \item \textbf{Beobachten:}
				Entwicklung liegt im erwarteten Rahmen.
  \item \textbf{Reagieren:}
				Grenzwert überschritten; Maßnahmen werden ausgelöst und terminiert.
  \item \textbf{Eskalieren:}
				kritischer Wert; Entscheidung über Pause/Korrektur/Abbruch und Reporting.
\end{itemize}

In Data-Science-Projekten sollten KRIs sowohl prozessuale als auch technische Signale abdecken.
Beispiele sind eine Verschlechterung der Datenqualität,
eine sinkende Modellgüte in kritischen Segmenten,
eine steigende Fehlerquote in der Datenpipeline
oder eine wachsende Abweichung zwischen Trainingsdaten und Produktionsdaten.
Solche Indikatoren eignen sich besonders,
um Drift und versteckte technische Abhängigkeiten früh sichtbar zu machen
\cite{Sculley2015}.

\subsection{Data Science spezifisches Monitoring im Sinne von MLOps}
Ein zentraler Risikotreiber ist die Veränderung der Modellleistung im Betrieb (z.B. Data/Concept Drift).
Neben Projektkennzahlen sind daher Daten- und Modellkennzahlen kontinuierlich zu überwachen \cite{Sculley2015,Amershi2019}.

Ein weiterer, für Data Science typischer Aspekt ist,
dass sich Risiken oft als technische Schuld in Datenabhängigkeiten und Modellabhängigkeiten
manifestieren.
Beispiele sind Abweichungen zwischen Training und Serving,
fragile Featuredefinitionen oder nicht dokumentierte Datenvorverarbeitung.
Monitoring sollte daher die Ende-zu-Ende-Kette von Datenquelle über Featureberechnung bis Inferenz abdecken \cite{Sculley2015}.

Ansätze aus dem MLOps Umfeld betonen,
dass Monitoring technisch unterstützt werden muss,
beispielsweise durch Logging, automatisierte Tests, Dashboards und Alerting,
damit Qualitätsabweichungen früh erkannt und reproduzierbar behoben werden können
\cite{Amershi2019}.
Damit wird Monitoring zum formalen Rückkopplungsmechanismus für Iterationen und Entscheidungspunkte im DASC PM \cite{schulz_dasc-pm_2022}.

Für das Monitoring in Data Science Projekten sind typischerweise folgende
Signalgruppen relevant \cite{Sculley2015,Amershi2019}:
\begin{itemize}
  \item \textbf{Datenqualität:}
				fehlende Werte, Ausreißer, Schemaänderungen, Dubletten,
				Verteilungsshifts in Eingangsmerkmalen.
  \item \textbf{Qualität von Labels und Ground Truth:}
				Verzögerungen in der Rückmeldung,
				Inkonsistenzen in der Annotation,
	  systematische Unterschiede zwischen Trainingsdaten und Realitätsdaten.
  \item \textbf{Modellgüte:}
				Abfall in Accuracy, F1 oder AUC,
				Anstieg von Fehlerquoten in kritischen Klassen.
  \item \textbf{Robustheit und Fairness:}
				Performanceunterschiede zwischen Gruppen,
				Stabilität gegenüber veränderten Datenlagen.
  \item \textbf{Betrieb und Sicherheit:}
				Latenzen, Ausfallraten, Zugriffsmuster,
				Änderungen an Datenpipelines oder Modellversionen.
  \item \textbf{Businesswirkung:}
				Nutzenkennzahlen wie Conversion, Kosten oder Durchlaufzeiten,
				sowie Anzeichen für Fehlsteuerung durch falsche Modellinterpretation
				im Geschäftsprozess.
\end{itemize}

\subsection{Einordnung über Standards und Managementsysteme}
Zur wissenschaftlichen Einordnung lässt sich das Monitoring zudem an
Normen für Managementsysteme für KI anknüpfen.
ISO und IEC fordern in ISO/IEC 42001 explizit,
dass KI Systeme über ihren Lebenszyklus hinweg überwacht
und Abweichungen kontrolliert behandelt werden \cite{ISOIEC42001_2023}.
Damit wird Risikomonitoring nicht nur als Projektpraxis,
sondern als Bestandteil organisationaler Steuerung und kontinuierlicher
Verbesserung verankert.

\subsection{Ergebnisartefakte}
Die Methode lässt sich durch wenige, klar definierte Artefakte operationalisieren.
Geeignete Ergebnisse sind:
\begin{itemize}
	\item ein \textbf{Monitoringplan} mit Taktung, Zuständigkeiten und Eskalationswegen
  \item ein \textbf{Katalog von Key Risk Indicators} mit Schwellenwerten und Reaktionen
  \item ein \textbf{aktualisiertes Risikeregister} mit Maßnahmenstatus und Entscheidungen
  \item ein \textbf{regelmäßiges Risikoreporting} an Stakeholder
  \item ein \textbf{Betriebsmonitoring} für Datenkennzahlen und Modellkennzahlen,
	umgesetzt über Logs, Dashboards und Alerts
	\cite{Amershi2019}
\end{itemize}
